\clearpage
\SourceLeaf{147}
\begin{center}
\textsc{Gauss} an \textsc{Olbers}. Göttingen, 24. Januar 1812.
\end{center}

\noindent . . . . . . . . . . . Hr.\ \textsc{Delambre} soll, wie ich höre, im franz.\ Moniteur eine ad modum suum sehr weitläuftige Chrie über die moindres carrés gegeben haben; ich lese den Moniteur nicht, welcher auch nur in jährlichen Lieferungen hierher kommt. Vielleicht finden Sie einmal Gelegenheit, \textls{öffentlich} zu bezeugen, dass ich Ihnen schon bei unserer ersten persönlichen Bekanntschaft im Jahr 1803 die Hauptmomente davon declarirt habe. Unter meinen Papieren finde ich, dass ich im Junius 1798, wo mir jene Methode eine \textls{längst} angewandte Sache war, zuerst \textsc{Laplaces} Methode gesehen und die Unverträglichkeit derselben mit den Grundsätzen der Wahrscheinlichkeitsrechnung in einem kurzen Notizen-Journal\HistoricalNote{[*)]}{[*)}{In dem auf Seite 20 bereits erwähnten Tagebuche von \textsc{Gauss} findet sich aus dem Jahre 1798 die Notiz: »\emph{Calculus probabilitatis contra La Place defensus. Gott. Jun. 17}«.]} über meine mathematischen Beschäftigungen angezeigt habe. Im Herbst 1802 habe ich die \textsc{viii}\textsuperscript{ten} Cereselemente in meinem Astronomischen Brouillonbuche nach der Methode der kleinsten Quadrate gefunden eingetragen. Die Papiere, worin ich in frühern Jahren, z.\,B.\ im Frühjahr 1799, auf \textsc{Ulugh Beighs} Zeitgleichungstafel jene Methode angewandt habe, sind verloren gegangen. Das einzige, worüber man sich wundern kann, ist, dass dieses Princip, was sich so leicht von selbst darbietet, dass man auf den Gedanken allein gar keinen besondern Werth legen kann, nicht schon 50 oder 100 Jahr früher von andern, z.\,B.\ \textsc{Euler} oder \textsc{Lambert} oder \textsc{Halley} oder \textsc{Tobias Mayer} angewandt ist, obwohl es ja sehr leicht sein kann, dass z.\,B.\ letzterer so etwas angewandt hat, ohne es zu proclamiren, so wie jeder Rechner nothwendig sich selbst eine Menge Vortheile und Methoden schafft, die er nur gelegentlich durch mündliche Tradition fortpflanzt.

\noindent . . . . . . . . . . . . .

